%%% BAT lecture 07 %%%
\documentclass{beamer}
\usepackage[utf8]{inputenc}
\usepackage{algorithm2e, amsmath, amssymb, amsfonts, graphicx}
% allow section.equation numbering
\numberwithin{equation}{section}
% use boadilla theme
\usetheme{Boadilla}
% remove navigation symbols
\usenavigationsymbolstemplate{}
% get numbered figure captions
\setbeamertemplate{caption}[numbered]
% changes itemize to circle + other things
\useoutertheme{split}
\useinnertheme{circles}

% command for the title string. change for each lecture
\newcommand{\lecturetitle}{Intro to Optimization, Part 2}
% allow automatic alert-highlighted references and hyperlinks
\newcommand{\aref}[1]{\alert{\ref{#1}}}
\newcommand{\ahref}[2]{\href{#1}{\alert{#2}}}
% title page stuff. brackets content displayed in footer bar
\title[\lecturetitle]{\lecturetitle}
% metadata. content in brackets is displayed in footer bar
\author[Derek Huang (BAC Advanced Team)]{Derek Huang}
\institute{BAC Advanced Team}
\date{May 21, 2021}

% change "ball" bullet to numbered bullet and section title for section
\setbeamertemplate{section in toc}{\inserttocsectionnumber.~\inserttocsection}
% change ball to gray square (copied from stackoverflow; \par needed for break)
\setbeamertemplate{subsection in toc}{        
    \hspace{1.2em}{\color{gray}\rule[0.3ex]{3pt}{3pt}}~\inserttocsubsection\par
}
% use default enumeration scheme
\setbeamertemplate{enumerate items}[default]
% required line that fixes the problem of \mathbf, \bf not working in beamer
% for later (post-2019) TeX Live installations. see the issue on GitHub:
% https://github.com/josephwright/beamer/issues/630
\DeclareFontShape{OT1}{cmss}{b}{n}{<->ssub * cmss/bx/n}{}

\begin{document}

% title slide
\begin{frame}
    \titlepage
    \centering
    % relative path may need to be updated depending on .tex file location
    \includegraphics[scale = 0.1]{../bac_logo1.png}
\end{frame}

% table of contents slide
\begin{frame}{Overview}
    \tableofcontents
\end{frame}

% section
\section{Unconstrained optimization}

% content slide
\begin{frame}{Motivation}
    \begin{itemize}
        \item
        Let $ \mathbf{X} \in \mathbb{R}^{N \times d} $ be the centered input
        matrix, $ \mathbf{y} \in \mathbb{R}^N $ the centered response
        vector\footnote{
            As a reminder, centering $ \mathbf{X} $, $ \mathbf{y} $ allows us
            to fit an interceptless model.
        }.
        Suppose we want to fit a lasso model to $ \mathbf{X}, \mathbf{y} $.

        \item
        In Lagrangian form, for $ \lambda \in (0, \infty) $, we must solve
        \begin{equation} \label{eq:lasso_obj}
            \begin{array}{ll}
                \displaystyle\min_\mathbf{w} &
                \Vert\mathbf{y} - \mathbf{Xw}\Vert_2^2 +
                \lambda\Vert\mathbf{w}\Vert_1
            \end{array}
        \end{equation}

        \item
        Theory tells us (\aref{eq:lasso_obj}) is convex but not
        differentiable. No closed form.

        \item
        Iterative methods exist\footnote{
            Coordinate descent and proximal gradient descent are often used.
        } to solve (\aref{eq:lasso_obj}) and similar
        problems. Software exists, but as users, we want to know \textbf{what}
        to use and \textbf{when}.

        \item
        \alert{We tailor our choice of algorithm to the problem at hand.}
    \end{itemize}
\end{frame}

\subsection{Line search methods}

\begin{frame}{Line search methods}
    \begin{itemize}
        \item
        Consider the generic unconstrained\footnote{
            Most machine learning problems can be written as unconstrained
            problems.
        } minimization problem
        \begin{equation} \label{eq:unconstrained_min}
            \begin{array}{ll}
                \displaystyle\min_\mathbf{x} & f(\mathbf{x})
            \end{array}
        \end{equation}
        $ f : \mathbb{R}^n \rightarrow \mathbb{R} $ is the objective. Let
        $ \mathbf{x}_{t - 1} \in \mathbb{R}^n $ denote the parameter
        estimate at the end of iteration $ t - 1 $, $ t \in \mathbb{N} $.

        \item
        Many algorithms for computing $ \mathbf{x}_t $ are
        \textit{line search methods}, which compute a \textit{search direction}
        $ \mathbf{d}_t \in \mathbb{R}^n $, \textit{step size}
        $ \eta_t \in (0, \infty) $ s.t. \cite{nocedal_opt}
        \begin{equation} \label{eq:line_search_eq}
            \mathbf{x}_t = \mathbf{x}_{t - 1} + \eta_t\mathbf{d}_t
        \end{equation}

        \item
        Obvious question is how to compute $ \mathbf{d}_t $, $ \eta_t $.
        Usually, $ \mathbf{d}_t $ is chosen to be a \textit{descent direction},
        i.e. so $ f(\mathbf{x}_{t - 1} + \eta_t\mathbf{d}_t) <
        f(\mathbf{x}_{t - 1}) $ \cite{nocedal_opt}. Note 
        $ \mathbf{d}_t^\top\nabla f(\mathbf{x}_{t - 1}) < 0 \Rightarrow $
        for small enough $ \eta_t $, $ \mathbf{d}_t $ is a descent direction
        \cite{nocedal_opt}.
    \end{itemize}

    % more spacing for the footnote
    \medskip
\end{frame}


% BibTeX slide for references. should use either acm or ieeetr style
\begin{frame}{References}
    \bibliographystyle{acm}
    % relative path may need to be updated depending on .tex file location
    \bibliography{../master_bib}
\end{frame}

\end{document}